\documentclass[11pt,a4paper]{article}

%====================================================
% PACKAGES
%====================================================
\usepackage[utf8]{inputenc}
\usepackage[T1]{fontenc}
\usepackage[french]{babel}
\usepackage[a4paper,margin=2cm]{geometry}
\usepackage{lmodern}
\usepackage{microtype}
\usepackage{amsmath,amssymb,amsthm,mathtools}
\usepackage{enumitem}
\usepackage{fancyhdr}
\usepackage{lastpage}
\usepackage{xcolor}
\usepackage[most]{tcolorbox}
\usepackage{hyperref}

%====================================================
% COULEURS
%====================================================
\definecolor{bleuprepa}{RGB}{20,55,110}
\definecolor{rougeprepa}{RGB}{155,35,35}
\definecolor{vertprepa}{RGB}{20,110,70}
\definecolor{orangeprepa}{RGB}{190,105,20}
\definecolor{grisclair}{RGB}{245,247,250}

\hypersetup{
  colorlinks=true,
  linkcolor=bleuprepa,
  urlcolor=bleuprepa
}

%====================================================
% MISE EN PAGE
%====================================================
\pagestyle{fancy}
\fancyhf{}
\fancyhead[L]{\textsc{CPGE scientifiques}}
\fancyhead[C]{\textsc{TD -- Fonctions réelles}}
\fancyhead[R]{\textsc{MPSI/PCSI/MP}}
\fancyfoot[C]{Page \thepage/\pageref{LastPage}}

\setlength{\parindent}{0pt}
\setlength{\parskip}{0.35em}

\setlist[enumerate]{label=\textbf{\arabic*.},leftmargin=1.1cm}
\setlist[itemize]{label=\(\triangleright\),leftmargin=1.2cm}

%====================================================
% BOÎTES
%====================================================
\tcbset{
  enhanced,
  breakable,
  sharp corners,
  boxrule=0.7pt,
  colback=white,
  fonttitle=\bfseries
}

\newcounter{exo}
\newcounter{corr}

\newcommand{\etoiles}[1]{\hfill{\small\textcolor{orangeprepa}{#1}}}
\newcommand{\R}{\mathbb{R}}
\newcommand{\N}{\mathbb{N}}
\newcommand{\eps}{\varepsilon}

\newtcolorbox{exercicebox}[2]{
  colframe=bleuprepa,
  colback=blue!2!white,
  title={Exercice #1 -- #2}
}

\newtcolorbox{correctionbox}[2]{
  colframe=vertprepa,
  colback=green!2!white,
  title={Correction #1 -- #2}
}

\newcommand{\exercice}[2]{
  \refstepcounter{exo}
  \begin{exercicebox}{\theexo}{#1}
  #2
  \end{exercicebox}
}

\newcommand{\correction}[2]{
  \refstepcounter{corr}
  \begin{correctionbox}{\thecorr}{#1}
  #2
  \end{correctionbox}
}

%====================================================
% DOCUMENT
%====================================================
\begin{document}

\begin{center}
{\Huge \bfseries TD complet -- Fonctions réelles}\\[0.2cm]
{\Large Limites, continuité, TVI, bijections et fonctions réciproques}\\[0.2cm]
{\large Classes préparatoires scientifiques -- MPSI, PCSI, MP}\\[0.25cm]
\textbf{Document prêt à compiler -- Énoncés et corrigés détaillés}
\end{center}

\vspace{0.3cm}

\begin{tcolorbox}[colframe=rougeprepa,colback=red!2!white,title={Objectifs du TD}]
Ce TD couvre les notions fondamentales du chapitre :
\begin{itemize}
  \item domaine de définition, image directe, image réciproque ;
  \item parité, périodicité, symétries ;
  \item limites finies, limites infinies, limites latérales ;
  \item caractérisation séquentielle des limites ;
  \item continuité, prolongement par continuité ;
  \item théorème des valeurs intermédiaires ;
  \item image d'un intervalle, image d'un segment ;
  \item injectivité, surjectivité, bijectivité ;
  \item théorème de la bijection continue et fonction réciproque.
\end{itemize}
\end{tcolorbox}

\tableofcontents
\newpage

%====================================================
\section{Exercices d'application directe -- APP}

\exercice{Domaines de définition simples \etoiles{★}}{
Déterminer le domaine de définition des fonctions suivantes :
\[
f_1(x)=\frac{1}{x-2},\qquad
f_2(x)=\sqrt{x+3},\qquad
f_3(x)=\ln(2x-1),
\]
\[
f_4(x)=\frac{\sqrt{x-1}}{x^2-4},\qquad
f_5(x)=\ln(x^2-1).
\]
Pour chaque fonction, rédiger les contraintes puis donner le domaine sous forme d'ensemble.
}

\exercice{Domaines de définition composés \etoiles{★★}}{
Déterminer le domaine de définition de :
\[
f(x)=\sqrt{\ln(x)},\qquad
g(x)=\ln\left(\frac{x+1}{x-2}\right),\qquad
h(x)=\frac{1}{\sqrt{4-x^2}}.
\]
On justifiera chaque contrainte imposée par les fonctions usuelles utilisées.
}

\exercice{Image directe d'un intervalle par une fonction quadratique \etoiles{★}}{
Soit
\[
f(x)=x^2-2x.
\]
\begin{enumerate}
  \item Écrire \(f(x)\) sous forme canonique.
  \item Déterminer \(f([0,3])\).
  \item Déterminer \(f([-2,0])\).
  \item Déterminer \(f([-1,4])\).
\end{enumerate}
}

\exercice{Image réciproque d'un intervalle \etoiles{★}}{
Soit
\[
f:\R\to\R,\qquad f(x)=x^2.
\]
Déterminer :
\[
f^{-1}([0,4]),\qquad f^{-1}([1,9]),\qquad f^{-1}(]-\infty,4[),\qquad f^{-1}(]4,+\infty[).
\]
On rappelle que cette notation désigne ici l'image réciproque d'un ensemble, et non la fonction réciproque.
}

\exercice{Parité \etoiles{★}}{
Étudier la parité des fonctions suivantes :
\[
f(x)=x^4-3x^2+1,\qquad
g(x)=\frac{x^3}{1+x^2},\qquad
h(x)=x^2+x.
\]
On commencera toujours par vérifier que le domaine est symétrique par rapport à \(0\).
}

\exercice{Périodicité \etoiles{★}}{
Montrer que les fonctions suivantes sont périodiques et déterminer une période :
\[
f(x)=\sin(3x),\qquad
g(x)=\cos(2x+\pi/4),\qquad
h(x)=\tan(5x).
\]
On ne demande pas nécessairement la période minimale, sauf si elle est immédiate.
}

\exercice{Limites de fonctions rationnelles \etoiles{★}}{
Calculer :
\[
\lim_{x\to 2}\frac{x^2-4}{x-2},\qquad
\lim_{x\to 1}\frac{x^3-1}{x-1},\qquad
\lim_{x\to -1}\frac{x^2-1}{x+1}.
\]
On factorisera avant de passer à la limite.
}

\exercice{Limites à l'infini de fonctions rationnelles \etoiles{★}}{
Calculer :
\[
\lim_{x\to+\infty}\frac{3x^2-x+1}{2x^2+5},\qquad
\lim_{x\to-\infty}\frac{x^3+1}{x^2+1},
\]
\[
\lim_{x\to+\infty}\frac{2x-1}{x^3+1}.
\]
On divisera par la puissance dominante.
}

\exercice{Limites latérales \etoiles{★★}}{
Étudier la limite en \(0\) de :
\[
f(x)=\frac{|x|}{x}.
\]
\begin{enumerate}
  \item Calculer \(f(x)\) pour \(x>0\).
  \item Calculer \(f(x)\) pour \(x<0\).
  \item Déterminer les limites à droite et à gauche en \(0\).
  \item Conclure sur l'existence de la limite en \(0\).
\end{enumerate}
}

\exercice{Gendarmes \etoiles{★★}}{
Montrer que :
\[
\lim_{x\to0}x\sin\left(\frac1x\right)=0,
\qquad
\lim_{x\to0}x^2\cos\left(\frac1{x^2}\right)=0.
\]
On précisera les encadrements utilisés.
}

%====================================================
\section{Exercices de méthode -- METH}

\exercice{Limite par la définition \(\eps,\eta\) \etoiles{★★}}{
Montrer directement avec la définition que :
\[
\lim_{x\to2}(3x+1)=7.
\]
\begin{enumerate}
  \item Calculer \(|(3x+1)-7|\).
  \item Trouver un choix possible de \(\eta\) en fonction de \(\eps\).
  \item Rédiger la preuve complète.
\end{enumerate}
}

\exercice{Limite quadratique avec la définition \etoiles{★★★}}{
Montrer avec la définition que :
\[
\lim_{x\to1}x^2=1.
\]
On pourra imposer d'abord \(|x-1|\leq1\) afin de contrôler \(|x+1|\).
}

\exercice{Limite infinie en un point \etoiles{★★}}{
Montrer avec la définition que :
\[
\lim_{x\to0}\frac1{x^2}=+\infty.
\]
\begin{enumerate}
  \item Rappeler la définition de \(f(x)\to+\infty\) lorsque \(x\to0\).
  \item Soit \(A>0\). Trouver une condition suffisante sur \(|x|\).
  \item Conclure.
\end{enumerate}
}

\exercice{Limite infinie à l'infini \etoiles{★★}}{
Montrer avec la définition que :
\[
\lim_{x\to+\infty}(x^2-3x)=+\infty.
\]
On pourra montrer que, pour \(x\geq6\),
\[
x^2-3x\geq \frac{x^2}{2}.
\]
}

\exercice{Caractérisation séquentielle : usage direct \etoiles{★★}}{
Soit \(f:D\to\R\), soit \(a\in\R\), et soit \(\ell\in\R\).
On suppose que :
\[
\lim_{x\to a}f(x)=\ell.
\]
\begin{enumerate}
  \item Soit \((x_n)\) une suite d'éléments de \(D\setminus\{a\}\) telle que \(x_n\to a\). Montrer que \(f(x_n)\to\ell\).
  \item Appliquer ce résultat à \(f(x)=x^2\), \(a=2\), et \(x_n=2+\frac1n\).
\end{enumerate}
}

\exercice{Caractérisation séquentielle : non-existence \etoiles{★★★}}{
Montrer que la fonction
\[
f(x)=\sin\left(\frac1x\right)
\]
n'admet pas de limite en \(0\).
On utilisera deux suites \(x_n\to0\) et \(y_n\to0\) donnant deux limites différentes pour \(f(x_n)\) et \(f(y_n)\).
}

\exercice{Continuité sur un domaine \etoiles{★★}}{
Étudier la continuité de la fonction
\[
f(x)=\frac{\sqrt{x+1}}{x^2-1}.
\]
\begin{enumerate}
  \item Déterminer son domaine de définition.
  \item Justifier que \(f\) est continue sur chaque intervalle de son domaine.
\end{enumerate}
}

\exercice{Continuité d'une fonction définie par morceaux \etoiles{★★}}{
Soit
\[
f(x)=
\begin{cases}
x^2+1, & x<1,\\
ax+b, & x\geq1.
\end{cases}
\]
\begin{enumerate}
  \item Pour quelles valeurs de \(a,b\) la fonction est-elle continue en \(1\) ?
  \item La continuité ailleurs pose-t-elle problème ?
  \item Donner deux exemples de couples \((a,b)\) convenables.
\end{enumerate}
}

\exercice{Prolongement par continuité \etoiles{★★}}{
Soit
\[
f(x)=\frac{x^2-1}{x-1}
\]
définie sur \(\R\setminus\{1\}\).
\begin{enumerate}
  \item Simplifier \(f(x)\) pour \(x\neq1\).
  \item Calculer \(\lim_{x\to1}f(x)\).
  \item Définir le prolongement continu de \(f\) en \(1\).
\end{enumerate}
}

\exercice{Prolongement avec racine \etoiles{★★★}}{
Soit
\[
f(x)=\frac{\sqrt{1+x}-1}{x}
\]
définie pour \(x>-1\) et \(x\neq0\).
\begin{enumerate}
  \item Multiplier par la quantité conjuguée.
  \item Calculer la limite de \(f(x)\) lorsque \(x\to0\).
  \item En déduire le prolongement par continuité en \(0\).
\end{enumerate}
}

%====================================================
\section{Exercices de démonstration -- DEM}

\exercice{Unicité de la limite \etoiles{★★}}{
Soit \(f:D\to\R\), et soit \(a\in\R\).
On suppose que :
\[
\lim_{x\to a}f(x)=\ell
\qquad\text{et}\qquad
\lim_{x\to a}f(x)=m.
\]
\begin{enumerate}
  \item Fixer \(\eps>0\), puis utiliser les deux définitions avec \(\eps/2\).
  \item Montrer que \(|\ell-m|\leq\eps\).
  \item Conclure que \(\ell=m\).
\end{enumerate}
}

\exercice{Limites latérales et limite bilatérale \etoiles{★★★}}{
Soit \(f\) définie au voisinage de \(a\), éventuellement privée de \(a\).
Montrer que :
\[
\lim_{x\to a}f(x)=\ell
\]
si et seulement si
\[
\lim_{x\to a^-}f(x)=\ell
\qquad\text{et}\qquad
\lim_{x\to a^+}f(x)=\ell.
\]
}

\exercice{Théorème des gendarmes pour les fonctions \etoiles{★★★}}{
Soient \(f,g,h\) trois fonctions définies au voisinage de \(a\), sauf éventuellement en \(a\).
On suppose que, au voisinage de \(a\),
\[
f(x)\leq g(x)\leq h(x),
\]
et que :
\[
\lim_{x\to a}f(x)=\ell,
\qquad
\lim_{x\to a}h(x)=\ell.
\]
Démontrer que :
\[
\lim_{x\to a}g(x)=\ell.
\]
}

\exercice{Passage à la limite dans une inégalité \etoiles{★★★}}{
Soient \(f,g\) deux fonctions définies au voisinage de \(a\), sauf éventuellement en \(a\), telles que
\[
f(x)\leq g(x)
\]
au voisinage de \(a\). On suppose que :
\[
f(x)\to\ell,\qquad g(x)\to m
\qquad (x\to a).
\]
Démontrer que :
\[
\ell\leq m.
\]
On pourra utiliser la caractérisation séquentielle.
}

\exercice{Continuité et opérations \etoiles{★★}}{
Soient \(f\) et \(g\) deux fonctions continues en \(a\).
\begin{enumerate}
  \item Montrer que \(f+g\) est continue en \(a\).
  \item Montrer que \(fg\) est continue en \(a\).
  \item Montrer que, si \(g(a)\neq0\), alors \(f/g\) est continue en \(a\).
\end{enumerate}
}

\exercice{Continuité d'une composée \etoiles{★★}}{
Soient \(g\) continue en \(a\) et \(f\) continue en \(g(a)\).
Démontrer que \(f\circ g\) est continue en \(a\).
}

\exercice{Théorème des valeurs intermédiaires : application abstraite \etoiles{★★★}}{
Soit \(f:[a,b]\to\R\) continue, avec \(f(a)f(b)<0\).
\begin{enumerate}
  \item Montrer que \(0\) est compris entre \(f(a)\) et \(f(b)\).
  \item En déduire qu'il existe \(c\in[a,b]\) tel que \(f(c)=0\).
  \item Expliquer pourquoi l'unicité n'est pas garantie.
\end{enumerate}
}

\exercice{Image d'un intervalle par une fonction continue \etoiles{★★★}}{
Soit \(I\) un intervalle de \(\R\), et soit \(f:I\to\R\) continue.
Démontrer que \(f(I)\) est un intervalle.
On utilisera le théorème des valeurs intermédiaires.
}

\exercice{Monotonie stricte et injectivité \etoiles{★★}}{
Soit \(I\) un intervalle de \(\R\), et soit \(f:I\to\R\) strictement croissante.
\begin{enumerate}
  \item Montrer que si \(x<y\), alors \(f(x)<f(y)\).
  \item En déduire que \(f\) est injective.
  \item Adapter la preuve au cas strictement décroissant.
\end{enumerate}
}

\exercice{Théorème de la bijection continue \etoiles{★★★}}{
Énoncer puis démontrer la partie suivante du théorème de la bijection continue :
si \(f:I\to\R\) est continue et strictement monotone, alors \(f\) réalise une bijection de \(I\) sur \(f(I)\).
On précisera où interviennent l'injectivité et la surjectivité.
}

%====================================================
\section{Exercices de réflexion -- PREP}

\exercice{Une fonction continue mais non injective \etoiles{★★}}{
Soit
\[
f:[-1,1]\to\R,\qquad f(x)=x^2.
\]
\begin{enumerate}
  \item Montrer que \(f\) est continue.
  \item Montrer que \(f\) n'est pas injective.
  \item Déterminer \(f([-1,1])\).
  \item Montrer que \(f\) devient bijective si on la restreint à \([0,1]\) avec un ensemble d'arrivée convenable.
\end{enumerate}
}

\exercice{TVI sans unicité \etoiles{★★}}{
Soit
\[
f(x)=x^2-1
\]
sur \([-2,2]\).
\begin{enumerate}
  \item Montrer que l'équation \(f(x)=0\) admet au moins une solution dans \([-2,2]\).
  \item Déterminer toutes les solutions.
  \item Expliquer pourquoi le TVI ne garantit pas l'unicité.
\end{enumerate}
}

\exercice{Existence et unicité sans dérivée \etoiles{★★★}}{
Montrer que l'équation
\[
x^3+x-1=0
\]
admet une unique solution réelle.
\begin{enumerate}
  \item Montrer l'existence d'une solution dans \([0,1]\).
  \item Montrer que la fonction \(x\mapsto x^3+x-1\) est strictement croissante sur \(\R\), sans utiliser la dérivée.
  \item Conclure.
\end{enumerate}
}

\exercice{Bijection de \(\R\) sur \(\R\) \etoiles{★★★}}{
Soit
\[
f(x)=x^3+x.
\]
\begin{enumerate}
  \item Montrer que \(f\) est continue sur \(\R\).
  \item Montrer que \(f\) est strictement croissante sur \(\R\), sans utiliser la dérivée.
  \item Déterminer les limites de \(f\) en \(\pm\infty\).
  \item Conclure que \(f\) réalise une bijection de \(\R\) sur \(\R\).
\end{enumerate}
}

\exercice{Bijection sur un intervalle ouvert \etoiles{★★★}}{
Soit
\[
f(x)=x+\frac1x
\]
définie sur \(]0,+\infty[\).
\begin{enumerate}
  \item Montrer que \(f(x)\geq2\) pour tout \(x>0\).
  \item Montrer que \(f\) n'est pas injective sur \(]0,+\infty[\).
  \item Montrer que \(f\) est injective sur \(]0,1]\) et sur \([1,+\infty[\).
  \item Déterminer l'image de chacun de ces deux intervalles.
\end{enumerate}
}

\exercice{Prolongement impossible \etoiles{★★}}{
Soit
\[
f(x)=\frac1x.
\]
\begin{enumerate}
  \item Calculer les limites à gauche et à droite en \(0\).
  \item Montrer que \(f\) n'est pas prolongeable par continuité en \(0\).
  \item Comparer avec la fonction \(x\mapsto \dfrac{\sin x}{x}\).
\end{enumerate}
}

\exercice{Fonction oscillante bornée \etoiles{★★★}}{
Soit
\[
f(x)=\sin\left(\frac1x\right)
\]
définie sur \(\R^*\).
\begin{enumerate}
  \item Montrer que \(f\) est bornée.
  \item Montrer que \(f\) n'admet pas de limite en \(0\).
  \item Montrer que \(x f(x)\) admet une limite en \(0\).
  \item Déterminer cette limite.
\end{enumerate}
}

\exercice{Image réciproque et fonction réciproque \etoiles{★★★}}{
Soit
\[
f:\R\to\R,\qquad f(x)=x^2.
\]
\begin{enumerate}
  \item Déterminer l'image réciproque \(f^{-1}(\{4\})\).
  \item La fonction \(f\) admet-elle une fonction réciproque de \(\R\) vers \(\R\) ?
  \item Restreindre \(f\) à \([0,+\infty[\) et déterminer sa réciproque.
  \item Expliquer la différence entre les deux notations \(f^{-1}(B)\) et \(f^{-1}(y)\).
\end{enumerate}
}

\exercice{Continuité et suites \etoiles{★★★}}{
Soit \(f:\R\to\R\) continue en \(a\).
\begin{enumerate}
  \item Montrer que pour toute suite \(x_n\to a\), on a \(f(x_n)\to f(a)\).
  \item Utiliser cette propriété pour montrer que la fonction indicatrice de \(\R_+\),
  \[
  \mathbf{1}_{\R_+}(x)=
  \begin{cases}
  1,&x\geq0,\\
  0,&x<0,
  \end{cases}
  \]
  n'est pas continue en \(0\).
\end{enumerate}
}

\exercice{Image d'un segment \etoiles{★★★}}{
Soit
\[
f(x)=x^2-4x+1
\]
sur \([0,5]\).
\begin{enumerate}
  \item Écrire \(f\) sous forme canonique.
  \item Déterminer le minimum de \(f\) sur \([0,5]\).
  \item Déterminer le maximum de \(f\) sur \([0,5]\).
  \item En déduire \(f([0,5])\).
\end{enumerate}
}

%====================================================
\section{Exercices type oraux -- Mines, Centrale, CCP}

\exercice{Oral CCP -- Équation \(x=\cos x\) \etoiles{★★★}}{
Montrer que l'équation
\[
x=\cos x
\]
admet une unique solution dans \([0,1]\).
\begin{enumerate}
  \item Poser une fonction adaptée.
  \item Montrer l'existence par le TVI.
  \item Montrer l'unicité sans utiliser la dérivée.
\end{enumerate}
}

\exercice{Oral CCP -- Fonction \(x+\ln x\) \etoiles{★★★}}{
Soit
\[
f(x)=x+\ln x
\]
définie sur \(]0,+\infty[\).
\begin{enumerate}
  \item Montrer que \(f\) est continue.
  \item Montrer que \(f\) est strictement croissante.
  \item Déterminer les limites aux bornes de l'intervalle.
  \item En déduire que \(f\) réalise une bijection de \(]0,+\infty[\) sur \(\R\).
\end{enumerate}
}

\exercice{Oral Mines -- Une fonction réciproque implicite \etoiles{★★★★}}{
Soit
\[
f(x)=x^3+x.
\]
\begin{enumerate}
  \item Montrer que \(f\) est bijective de \(\R\) sur \(\R\).
  \item On note \(g=f^{-1}\). Déterminer \(g(0)\).
  \item Montrer que \(g\) est impaire.
  \item Résoudre \(g(y)=1\).
\end{enumerate}
}

\exercice{Oral Centrale -- Paramètre et TVI \etoiles{★★★★}}{
Soit \(a\in\R\). On considère l'équation
\[
x^3+x=a.
\]
\begin{enumerate}
  \item Montrer que, pour tout \(a\in\R\), elle admet une unique solution réelle.
  \item On note cette solution \(\varphi(a)\). Montrer que \(\varphi\) est une fonction de \(\R\) vers \(\R\).
  \item Interpréter \(\varphi\) comme fonction réciproque d'une fonction usuelle de ce TD.
\end{enumerate}
}

\exercice{Oral Mines -- Fonction définie par morceaux \etoiles{★★★★}}{
Soit
\[
f(x)=
\begin{cases}
\sin(1/x),&x\neq0,\\
0,&x=0.
\end{cases}
\]
\begin{enumerate}
  \item Montrer que \(f\) est continue sur \(\R^*\).
  \item Montrer que \(f\) n'est pas continue en \(0\).
  \item Étudier la continuité en \(0\) de \(g(x)=x\sin(1/x)\), avec \(g(0)=0\).
\end{enumerate}
}

\exercice{Oral Centrale -- Prolongement et paramètre \etoiles{★★★★}}{
Pour \(a,b\in\R\), on définit
\[
f(x)=
\begin{cases}
\dfrac{\sqrt{1+x}-1}{x},&x\neq0,\ x>-1,\\
a,&x=0,
\end{cases}
\]
puis
\[
g(x)=bf(x)+x.
\]
\begin{enumerate}
  \item Déterminer \(a\) pour que \(f\) soit continue en \(0\).
  \item Pour cette valeur de \(a\), étudier la continuité de \(g\) en \(0\).
  \item Déterminer \(g(0)\).
\end{enumerate}
}

\exercice{Oral CCP -- Image d'un intervalle \etoiles{★★★}}{
Soit \(f:[0,2]\to\R\) définie par
\[
f(x)=x^2-2x+3.
\]
\begin{enumerate}
  \item Montrer que \(f\) est continue.
  \item Déterminer son image sur \([0,2]\).
  \item Montrer directement que cette image est un segment.
\end{enumerate}
}

\exercice{Oral Mines -- Non-existence par suites \etoiles{★★★★}}{
Étudier l'existence de la limite en \(0\) de
\[
f(x)=\cos\left(\frac1x\right).
\]
On proposera deux suites tendant vers \(0\) et donnant deux limites différentes.
}

\exercice{Oral Centrale -- Bijection trigonométrique \etoiles{★★★★}}{
Soit
\[
f:\left[-\frac{\pi}{2},\frac{\pi}{2}\right]\to[-1,1],
\qquad f(x)=\sin x.
\]
\begin{enumerate}
  \item Justifier que \(f\) est continue.
  \item Justifier que \(f\) est strictement croissante sur cet intervalle.
  \item Montrer que \(f\) est bijective.
  \item Nommer sa fonction réciproque.
\end{enumerate}
}

\exercice{Oral Mines -- Fonction exponentielle \etoiles{★★★}}{
Montrer que
\[
\exp:\R\to]0,+\infty[
\]
est une bijection continue strictement croissante. Identifier sa fonction réciproque et donner ses propriétés de continuité et de monotonie.
}

%====================================================
\section{Problèmes longs type concours}

\exercice{Problème long 1 -- Étude complète d'une bijection \etoiles{★★★★★}}{
On considère
\[
f:\R\to\R,\qquad f(x)=x^3+x.
\]
\begin{enumerate}
  \item Montrer que \(f\) est continue sur \(\R\).
  \item Montrer que \(f\) est strictement croissante sans utiliser la dérivée.
  \item Déterminer \(\lim_{x\to+\infty}f(x)\) et \(\lim_{x\to-\infty}f(x)\).
  \item Montrer que \(f\) réalise une bijection de \(\R\) sur \(\R\).
  \item On note \(g=f^{-1}\). Montrer que \(g\) est continue et strictement croissante.
  \item Montrer que \(g\) est impaire.
  \item Résoudre \(g(y)=2\).
\end{enumerate}
}

\exercice{Problème long 2 -- TVI, unicité et paramètre \etoiles{★★★★★}}{
Pour \(a\in\R\), on considère l'équation
\[
(E_a)\qquad x^5+x=a.
\]
\begin{enumerate}
  \item Montrer que, pour tout \(a\in\R\), l'équation \((E_a)\) admet au moins une solution réelle.
  \item Montrer que cette solution est unique.
  \item On note \(\varphi(a)\) cette unique solution. Définir rigoureusement \(\varphi\).
  \item Montrer que \(\varphi\) est la réciproque d'une fonction continue strictement croissante.
  \item En déduire que \(\varphi\) est continue et strictement croissante.
\end{enumerate}
}

\exercice{Problème long 3 -- Prolongement et continuité \etoiles{★★★★}}{
Soit
\[
f(x)=\frac{\sqrt{1+x}-1}{x}
\]
définie sur \(]-1,+\infty[\setminus\{0\}\).
\begin{enumerate}
  \item Déterminer le domaine de définition.
  \item Simplifier l'expression à l'aide de la quantité conjuguée.
  \item Calculer la limite en \(0\).
  \item Définir un prolongement \(\widetilde f\) continu en \(0\).
  \item Étudier la continuité de \(\widetilde f\) sur \(]-1,+\infty[\).
\end{enumerate}
}

\exercice{Problème long 4 -- Fonction réciproque sur un intervalle \etoiles{★★★★}}{
Soit
\[
f:[0,+\infty[\to[1,+\infty[,\qquad f(x)=x^2+1.
\]
\begin{enumerate}
  \item Montrer que \(f\) est continue.
  \item Montrer que \(f\) est strictement croissante.
  \item Déterminer \(f([0,+\infty[)\).
  \item Montrer que \(f\) est bijective de \([0,+\infty[\) sur \([1,+\infty[\).
  \item Déterminer explicitement \(f^{-1}\).
  \item Vérifier les identités \(f^{-1}\circ f=\operatorname{Id}_{[0,+\infty[}\) et \(f\circ f^{-1}=\operatorname{Id}_{[1,+\infty[}\).
\end{enumerate}
}

\exercice{Problème long 5 -- Limites latérales, continuité et recollement \etoiles{★★★★}}{
Soient \(a,b\in\R\). On définit
\[
f(x)=
\begin{cases}
x^2+a,&x<1,\\
bx+1,&x\geq1.
\end{cases}
\]
\begin{enumerate}
  \item Étudier la continuité de \(f\) sur \(]-\infty,1[\) et sur \(]1,+\infty[\).
  \item Calculer la limite à gauche en \(1\).
  \item Calculer la limite à droite en \(1\).
  \item Donner une condition nécessaire et suffisante sur \(a,b\) pour que \(f\) soit continue en \(1\).
  \item Donner trois couples \((a,b)\) convenables.
\end{enumerate}
}

%====================================================
\newpage
\section{Corrigé complet}

\correction{Domaines de définition simples}{
\[
D_{f_1}=\R\setminus\{2\}.
\]
Pour \(f_2\), il faut \(x+3\geq0\), donc
\[
D_{f_2}=[-3,+\infty[.
\]
Pour \(f_3\), il faut \(2x-1>0\), donc
\[
D_{f_3}=]1/2,+\infty[.
\]
Pour \(f_4\), il faut \(x-1\geq0\) et \(x^2-4\neq0\). Donc \(x\geq1\) et \(x\neq2\). Ainsi
\[
D_{f_4}=[1,+\infty[\setminus\{2\}.
\]
Pour \(f_5\), il faut \(x^2-1>0\), donc \(|x|>1\). Ainsi
\[
D_{f_5}=]-\infty,-1[\cup]1,+\infty[.
\]
}

\correction{Domaines de définition composés}{
Pour \(f(x)=\sqrt{\ln x}\), il faut \(x>0\) et \(\ln x\geq0\). Donc \(x\geq1\), d'où
\[
D_f=[1,+\infty[.
\]
Pour \(g\), il faut
\[
\frac{x+1}{x-2}>0,\qquad x\neq2.
\]
Un tableau de signes donne
\[
D_g=]-\infty,-1[\cup]2,+\infty[.
\]
Pour \(h\), il faut
\[
4-x^2>0.
\]
Donc
\[
D_h=]-2,2[.
\]
}

\correction{Image directe d'un intervalle par une fonction quadratique}{
On écrit
\[
f(x)=x^2-2x=(x-1)^2-1.
\]
Sur \([0,3]\), le minimum vaut \(-1\) en \(1\), et le maximum vaut \(3\) en \(3\). Donc
\[
f([0,3])=[-1,3].
\]
Sur \([-2,0]\), la fonction est décroissante car l'intervalle est à gauche de \(1\). On obtient
\[
f(-2)=8,\qquad f(0)=0,
\]
donc
\[
f([-2,0])=[0,8].
\]
Sur \([-1,4]\), le minimum vaut \(-1\) en \(1\), et
\[
f(-1)=3,\qquad f(4)=8.
\]
Donc
\[
f([-1,4])=[-1,8].
\]
}

\correction{Image réciproque d'un intervalle}{
On utilise \(f^{-1}(B)=\{x\in\R:x^2\in B\}\).
\[
f^{-1}([0,4])=[-2,2].
\]
\[
f^{-1}([1,9])=[-3,-1]\cup[1,3].
\]
\[
f^{-1}(]-\infty,4[)=]-2,2[.
\]
\[
f^{-1}(]4,+\infty[)=]-\infty,-2[\cup]2,+\infty[.
\]
}

\correction{Parité}{
Les trois fonctions sont définies sur \(\R\), qui est symétrique par rapport à \(0\).
\[
f(-x)=(-x)^4-3(-x)^2+1=x^4-3x^2+1=f(x),
\]
donc \(f\) est paire.
\[
g(-x)=\frac{(-x)^3}{1+(-x)^2}=-\frac{x^3}{1+x^2}=-g(x),
\]
donc \(g\) est impaire.
Enfin
\[
h(-x)=x^2-x.
\]
Ce n'est ni \(h(x)=x^2+x\), ni \(-h(x)=-x^2-x\). Donc \(h\) n'est ni paire ni impaire.
}

\correction{Périodicité}{
Pour \(f(x)=\sin(3x)\), on cherche \(T>0\) tel que \(3T=2\pi\), donc
\[
T=\frac{2\pi}{3}.
\]
Pour \(g(x)=\cos(2x+\pi/4)\), il suffit que \(2T=2\pi\), donc
\[
T=\pi.
\]
Pour \(h(x)=\tan(5x)\), la tangente est de période \(\pi\). Il suffit que \(5T=\pi\), donc
\[
T=\frac{\pi}{5}.
\]
}

\correction{Limites de fonctions rationnelles}{
Pour \(x\neq2\),
\[
\frac{x^2-4}{x-2}=x+2,
\]
donc la limite vaut \(4\).
Pour \(x\neq1\),
\[
\frac{x^3-1}{x-1}=x^2+x+1,
\]
donc la limite vaut \(3\).
Pour \(x\neq-1\),
\[
\frac{x^2-1}{x+1}=x-1,
\]
donc la limite vaut \(-2\).
}

\correction{Limites à l'infini de fonctions rationnelles}{
On divise par la puissance dominante.
\[
\frac{3x^2-x+1}{2x^2+5}
=
\frac{3-\frac1x+\frac1{x^2}}{2+\frac5{x^2}}
\to\frac32.
\]
\[
\frac{x^3+1}{x^2+1}
=
x\frac{1+\frac1{x^3}}{1+\frac1{x^2}}.
\]
Lorsque \(x\to-\infty\), cette expression tend vers \(-\infty\).
Enfin
\[
\frac{2x-1}{x^3+1}
=
\frac{\frac2{x^2}-\frac1{x^3}}{1+\frac1{x^3}}
\to0.
\]
}

\correction{Limites latérales}{
Si \(x>0\), alors \(|x|=x\), donc
\[
f(x)=1.
\]
Si \(x<0\), alors \(|x|=-x\), donc
\[
f(x)=-1.
\]
Ainsi
\[
\lim_{x\to0^+}f(x)=1,\qquad
\lim_{x\to0^-}f(x)=-1.
\]
Les limites latérales sont différentes. La limite en \(0\) n'existe donc pas.
}

\correction{Gendarmes}{
Pour \(x\neq0\),
\[
-1\leq\sin(1/x)\leq1.
\]
Donc
\[
-|x|\leq x\sin(1/x)\leq |x|.
\]
Comme \(|x|\to0\), le théorème des gendarmes donne
\[
x\sin(1/x)\to0.
\]
De même,
\[
-1\leq\cos(1/x^2)\leq1,
\]
donc
\[
-x^2\leq x^2\cos(1/x^2)\leq x^2.
\]
Les deux bornes tendent vers \(0\). Donc
\[
x^2\cos(1/x^2)\to0.
\]
}

\correction{Limite par la définition \(\eps,\eta\)}{
On calcule :
\[
|(3x+1)-7|=|3x-6|=3|x-2|.
\]
Soit \(\eps>0\). On choisit
\[
\eta=\frac{\eps}{3}.
\]
Alors, si \(|x-2|\leq\eta\),
\[
|(3x+1)-7|=3|x-2|\leq3\eta=\eps.
\]
Donc
\[
\lim_{x\to2}(3x+1)=7.
\]
}

\correction{Limite quadratique avec la définition}{
On veut contrôler
\[
|x^2-1|=|x-1||x+1|.
\]
Imposons \(|x-1|\leq1\). Alors \(0\leq x\leq2\), donc \(|x+1|\leq3\).
Ainsi
\[
|x^2-1|\leq3|x-1|.
\]
Soit \(\eps>0\). On choisit
\[
\eta=\min\left(1,\frac{\eps}{3}\right).
\]
Alors, si \(|x-1|\leq\eta\),
\[
|x^2-1|\leq3|x-1|\leq3\eta\leq\eps.
\]
Donc \(x^2\to1\) lorsque \(x\to1\).
}

\correction{Limite infinie en un point}{
Soit \(A>0\). On veut
\[
\frac1{x^2}\geq A.
\]
Cela équivaut à
\[
x^2\leq \frac1A,
\]
donc
\[
|x|\leq \frac1{\sqrt A}.
\]
On choisit
\[
\eta=\frac1{\sqrt A}.
\]
Alors, si \(0<|x|\leq\eta\), on a
\[
\frac1{x^2}\geq A.
\]
Donc
\[
\lim_{x\to0}\frac1{x^2}=+\infty.
\]
}

\correction{Limite infinie à l'infini}{
Pour \(x\geq6\),
\[
x^2-3x=x(x-3)\geq \frac{x^2}{2}.
\]
Soit \(B\in\R\). Si \(B\leq0\), tout \(x\) assez grand convient. Si \(B>0\), on choisit
\[
A=\max(6,\sqrt{2B}).
\]
Alors, pour \(x\geq A\),
\[
x^2-3x\geq\frac{x^2}{2}\geq B.
\]
Donc
\[
x^2-3x\to+\infty.
\]
}

\correction{Caractérisation séquentielle : usage direct}{
Par caractérisation séquentielle de la limite, si \(x_n\to a\) et \(x_n\neq a\), alors
\[
f(x_n)\to\ell.
\]
Pour \(f(x)=x^2\), \(a=2\), \(x_n=2+1/n\), on a \(x_n\to2\). Donc
\[
f(x_n)=\left(2+\frac1n\right)^2\to4.
\]
}

\correction{Caractérisation séquentielle : non-existence}{
Prenons
\[
x_n=\frac{1}{\frac{\pi}{2}+2n\pi},
\qquad
y_n=\frac{1}{-\frac{\pi}{2}+2n\pi}
\]
pour \(n\) assez grand. Alors \(x_n\to0\) et \(y_n\to0\). De plus,
\[
f(x_n)=\sin\left(\frac{\pi}{2}+2n\pi\right)=1,
\]
et
\[
f(y_n)=\sin\left(-\frac{\pi}{2}+2n\pi\right)=-1.
\]
Les deux limites sont distinctes. Donc \(f\) n'admet pas de limite en \(0\).
}

\correction{Continuité sur un domaine}{
Il faut \(x+1\geq0\), donc \(x\geq-1\), et \(x^2-1\neq0\), donc \(x\neq-1\) et \(x\neq1\).
Ainsi
\[
D_f=]-1,+\infty[\setminus\{1\}.
\]
Sur chacun des intervalles \(]-1,1[\) et \(]1,+\infty[\), la fonction est un quotient de fonctions continues dont le dénominateur ne s'annule pas. Elle est donc continue sur ces deux intervalles.
}

\correction{Continuité d'une fonction définie par morceaux}{
À gauche de \(1\),
\[
\lim_{x\to1^-}f(x)=1^2+1=2.
\]
À droite,
\[
\lim_{x\to1^+}f(x)=a+b.
\]
Comme \(f(1)=a+b\), la continuité en \(1\) équivaut à
\[
a+b=2.
\]
Ailleurs, chaque morceau est polynomial, donc continu. Des exemples :
\[
(a,b)=(1,1),\quad (2,0),\quad (0,2).
\]
}

\correction{Prolongement par continuité}{
Pour \(x\neq1\),
\[
\frac{x^2-1}{x-1}=\frac{(x-1)(x+1)}{x-1}=x+1.
\]
Donc
\[
\lim_{x\to1}f(x)=2.
\]
Le prolongement continu \(\widetilde f\) est défini par
\[
\widetilde f(x)=f(x)\quad(x\neq1),
\qquad
\widetilde f(1)=2.
\]
}

\correction{Prolongement avec racine}{
Pour \(x\neq0\),
\[
f(x)=\frac{\sqrt{1+x}-1}{x}\cdot
\frac{\sqrt{1+x}+1}{\sqrt{1+x}+1}
=
\frac{x}{x(\sqrt{1+x}+1)}
=
\frac1{\sqrt{1+x}+1}.
\]
Donc
\[
\lim_{x\to0}f(x)=\frac12.
\]
On prolonge par continuité en posant
\[
\widetilde f(0)=\frac12.
\]
}

\correction{Unicité de la limite}{
Soit \(\eps>0\). Puisque \(f(x)\to\ell\), il existe \(\eta_1>0\) tel que
\[
0<|x-a|\leq\eta_1\Rightarrow |f(x)-\ell|\leq\frac{\eps}{2}.
\]
Puisque \(f(x)\to m\), il existe \(\eta_2>0\) tel que
\[
0<|x-a|\leq\eta_2\Rightarrow |f(x)-m|\leq\frac{\eps}{2}.
\]
Pour \(0<|x-a|\leq\min(\eta_1,\eta_2)\),
\[
|\ell-m|\leq|\ell-f(x)|+|f(x)-m|\leq\eps.
\]
Ainsi, pour tout \(\eps>0\), \(|\ell-m|\leq\eps\). Donc \(\ell=m\).
}

\correction{Limites latérales et limite bilatérale}{
Si \(f(x)\to\ell\) lorsque \(x\to a\), alors la propriété reste vraie en se restreignant aux \(x<a\) et aux \(x>a\). Les limites latérales valent donc \(\ell\).

Réciproquement, supposons les deux limites latérales égales à \(\ell\). Pour \(\eps>0\), il existe \(\eta_1>0\) tel que
\[
a-\eta_1\leq x<a\Rightarrow |f(x)-\ell|\leq\eps,
\]
et \(\eta_2>0\) tel que
\[
a<x\leq a+\eta_2\Rightarrow |f(x)-\ell|\leq\eps.
\]
Avec \(\eta=\min(\eta_1,\eta_2)\), si \(0<|x-a|\leq\eta\), alors \(x<a\) ou \(x>a\), et dans les deux cas
\[
|f(x)-\ell|\leq\eps.
\]
Donc \(f(x)\to\ell\).
}

\correction{Théorème des gendarmes pour les fonctions}{
Soit \(\eps>0\). Comme \(f(x)\to\ell\), on a, pour \(x\) assez proche de \(a\),
\[
f(x)\geq\ell-\eps.
\]
Comme \(h(x)\to\ell\), on a, pour \(x\) assez proche de \(a\),
\[
h(x)\leq\ell+\eps.
\]
Au voisinage de \(a\),
\[
f(x)\leq g(x)\leq h(x).
\]
Donc, pour \(x\) assez proche de \(a\),
\[
\ell-\eps\leq g(x)\leq\ell+\eps.
\]
Ainsi
\[
|g(x)-\ell|\leq\eps.
\]
Donc \(g(x)\to\ell\).
}

\correction{Passage à la limite dans une inégalité}{
Soit \((x_n)\) une suite tendant vers \(a\), avec \(x_n\neq a\). Pour \(n\) assez grand,
\[
f(x_n)\leq g(x_n).
\]
Par caractérisation séquentielle,
\[
f(x_n)\to\ell,\qquad g(x_n)\to m.
\]
En passant à la limite dans l'inégalité entre suites, on obtient
\[
\ell\leq m.
\]
}

\correction{Continuité et opérations}{
Puisque \(f\) et \(g\) sont continues en \(a\),
\[
f(x)\to f(a),\qquad g(x)\to g(a).
\]
Par opérations sur les limites,
\[
f(x)+g(x)\to f(a)+g(a),
\]
donc \(f+g\) est continue en \(a\).
De même,
\[
f(x)g(x)\to f(a)g(a),
\]
donc \(fg\) est continue en \(a\).
Si \(g(a)\neq0\), alors \(g(x)\neq0\) au voisinage de \(a\), et
\[
\frac{f(x)}{g(x)}\to\frac{f(a)}{g(a)}.
\]
Donc \(f/g\) est continue en \(a\).
}

\correction{Continuité d'une composée}{
Comme \(g\) est continue en \(a\),
\[
g(x)\to g(a).
\]
Comme \(f\) est continue en \(g(a)\),
\[
f(y)\to f(g(a))\quad\text{lorsque }y\to g(a).
\]
En prenant \(y=g(x)\), on obtient
\[
f(g(x))\to f(g(a)).
\]
Donc \(f\circ g\) est continue en \(a\).
}

\correction{Théorème des valeurs intermédiaires : application abstraite}{
Si \(f(a)f(b)<0\), alors \(f(a)\) et \(f(b)\) sont de signes opposés. Ainsi \(0\) est compris entre \(f(a)\) et \(f(b)\). Comme \(f\) est continue sur \([a,b]\), le TVI assure qu'il existe \(c\in[a,b]\) tel que
\[
f(c)=0.
\]
Le TVI ne donne pas l'unicité : une fonction continue peut couper plusieurs fois l'axe des abscisses.
}

\correction{Image d'un intervalle par une fonction continue}{
Soient \(u,v\in f(I)\), avec \(u\leq v\). Il existe \(a,b\in I\) tels que
\[
u=f(a),\qquad v=f(b).
\]
Soit \(y\in[u,v]\). Comme \(I\) est un intervalle, le segment entre \(a\) et \(b\) est inclus dans \(I\). La fonction \(f\) est continue sur ce segment. Par le TVI, il existe \(c\) entre \(a\) et \(b\) tel que
\[
f(c)=y.
\]
Donc \(y\in f(I)\). Ainsi \(f(I)\) contient tout réel compris entre deux de ses éléments : c'est un intervalle.
}

\correction{Monotonie stricte et injectivité}{
Si \(f\) est strictement croissante et si \(x<y\), alors
\[
f(x)<f(y).
\]
Soient \(x,y\in I\) tels que \(f(x)=f(y)\). Si \(x<y\), on aurait \(f(x)<f(y)\), contradiction. Si \(y<x\), on aurait \(f(y)<f(x)\), contradiction. Donc \(x=y\). Ainsi \(f\) est injective.
Le cas strictement décroissant se traite de même.
}

\correction{Théorème de la bijection continue}{
Si \(f:I\to\R\) est continue et strictement monotone, alors \(f\) est injective par stricte monotonie. Par ailleurs, une fonction est toujours surjective de son ensemble de départ vers son image. Donc
\[
f:I\to f(I)
\]
est surjective. Elle est donc bijective de \(I\) sur \(f(I)\). La continuité permet en outre de savoir que \(f(I)\) est un intervalle, ce qui donne la forme analytique de l'ensemble d'arrivée naturel.
}

\correction{Une fonction continue mais non injective}{
La fonction \(x\mapsto x^2\) est polynomiale, donc continue. Elle n'est pas injective sur \([-1,1]\), car
\[
f(-1)=f(1)=1
\]
avec \(-1\neq1\). Sur \([-1,1]\),
\[
f([-1,1])=[0,1].
\]
Restreinte à \([0,1]\), la fonction est strictement croissante et continue, donc bijective de \([0,1]\) sur \([0,1]\).
}

\correction{TVI sans unicité}{
La fonction \(f(x)=x^2-1\) est continue sur \([-2,2]\). On a
\[
f(0)=-1,\qquad f(2)=3.
\]
Le TVI donne au moins une solution à \(f(x)=0\). En réalité,
\[
x^2-1=0
\Longleftrightarrow
x=\pm1.
\]
Il y a donc deux solutions. Le TVI garantit une existence, pas une unicité.
}

\correction{Existence et unicité sans dérivée}{
Posons
\[
f(x)=x^3+x-1.
\]
On a
\[
f(0)=-1,\qquad f(1)=1.
\]
Comme \(f\) est continue, le TVI donne une solution dans \([0,1]\).

Si \(x<y\), alors
\[
f(y)-f(x)=y^3-x^3+y-x
=(y-x)(y^2+xy+x^2+1).
\]
Or
\[
y^2+xy+x^2=\frac12(x+y)^2+\frac12(x^2+y^2)\geq0.
\]
Donc \(y^2+xy+x^2+1>0\). Comme \(y-x>0\),
\[
f(y)-f(x)>0.
\]
Donc \(f\) est strictement croissante, donc injective. L'équation admet donc une unique solution réelle.
}

\correction{Bijection de \(\R\) sur \(\R\)}{
La fonction \(f(x)=x^3+x\) est polynomiale, donc continue sur \(\R\). Si \(x<y\),
\[
f(y)-f(x)=y^3-x^3+y-x
=(y-x)(y^2+xy+x^2+1)>0.
\]
Donc \(f\) est strictement croissante. De plus,
\[
\lim_{x\to+\infty}(x^3+x)=+\infty,\qquad
\lim_{x\to-\infty}(x^3+x)=-\infty.
\]
Par le théorème de la bijection continue, \(f\) réalise une bijection de \(\R\) sur \(\R\).
}

\correction{Bijection sur un intervalle ouvert}{
Pour \(x>0\),
\[
x+\frac1x\geq2
\]
par l'inégalité classique \(t+1/t\geq2\). La fonction n'est pas injective sur \(]0,+\infty[\), car
\[
f(2)=2+\frac12=\frac52=f\left(\frac12\right).
\]
Sur \(]0,1]\), si \(0<x<y\leq1\), alors \(1/x>1/y\), et on vérifie que \(f\) décroît strictement. Sur \([1,+\infty[\), \(f\) croît strictement.
On obtient
\[
f(]0,1])=[2,+\infty[,
\qquad
f([1,+\infty[)=[2,+\infty[.
\]
}

\correction{Prolongement impossible}{
Lorsque \(x\to0^+\),
\[
\frac1x\to+\infty.
\]
Lorsque \(x\to0^-\),
\[
\frac1x\to-\infty.
\]
Il n'existe donc pas de limite finie en \(0\). La fonction n'est pas prolongeable par continuité.
En revanche,
\[
\frac{\sin x}{x}\to1
\]
lorsque \(x\to0\), donc cette dernière fonction est prolongeable en posant sa valeur en \(0\) égale à \(1\).
}

\correction{Fonction oscillante bornée}{
On a toujours
\[
-1\leq \sin(1/x)\leq1,
\]
donc \(f\) est bornée. Elle n'a pas de limite en \(0\), comme montré par deux suites donnant \(1\) et \(-1\). En revanche,
\[
-|x|\leq x\sin(1/x)\leq |x|.
\]
Par encadrement,
\[
x\sin(1/x)\to0.
\]
}

\correction{Image réciproque et fonction réciproque}{
\[
f^{-1}(\{4\})=\{-2,2\}.
\]
La fonction \(x\mapsto x^2\) n'est pas injective sur \(\R\), donc elle n'admet pas de fonction réciproque de \(\R\) vers \(\R\).
Restreinte à \([0,+\infty[\), elle est bijective de \([0,+\infty[\) vers \([0,+\infty[\), et sa réciproque est
\[
y\mapsto\sqrt y.
\]
L'image réciproque \(f^{-1}(B)\) existe pour tout ensemble \(B\). La fonction réciproque \(f^{-1}(y)\) n'existe que si \(f\) est bijective.
}

\correction{Continuité et suites}{
Si \(f\) est continue en \(a\), alors, pour toute suite \(x_n\to a\), la caractérisation séquentielle donne
\[
f(x_n)\to f(a).
\]
Pour \(\mathbf{1}_{\R_+}\), prenons
\[
x_n=-\frac1n,\qquad y_n=\frac1n.
\]
Alors \(x_n\to0\), \(y_n\to0\), mais
\[
\mathbf{1}_{\R_+}(x_n)=0,\qquad
\mathbf{1}_{\R_+}(y_n)=1.
\]
La fonction n'est donc pas continue en \(0\).
}

\correction{Image d'un segment}{
\[
f(x)=x^2-4x+1=(x-2)^2-3.
\]
Le minimum sur \([0,5]\) est donc \(-3\), atteint en \(2\).
Aux bornes :
\[
f(0)=1,\qquad f(5)=25-20+1=6.
\]
Le maximum vaut donc \(6\), atteint en \(5\). Comme \(f\) est continue sur le segment,
\[
f([0,5])=[-3,6].
\]
}

\correction{Oral CCP -- Équation \(x=\cos x\)}{
Posons
\[
f(x)=\cos x-x.
\]
La fonction est continue sur \([0,1]\). On a
\[
f(0)=1>0,
\qquad
f(1)=\cos1-1<0.
\]
Par le TVI, il existe \(c\in[0,1]\) tel que \(f(c)=0\), donc \(c=\cos c\).

Pour l'unicité, si \(0\leq x<y\leq1\), alors \(\cos x>\cos y\) sur \([0,1]\), donc
\[
\cos x-x>\cos y-y.
\]
Ainsi \(f\) est strictement décroissante, donc injective. Il y a une unique solution.
}

\correction{Oral CCP -- Fonction \(x+\ln x\)}{
La fonction \(x\mapsto x+\ln x\) est continue sur \(]0,+\infty[\). Elle est strictement croissante comme somme de deux fonctions strictement croissantes.
De plus,
\[
\lim_{x\to0^+}(x+\ln x)=-\infty,
\qquad
\lim_{x\to+\infty}(x+\ln x)=+\infty.
\]
Par le théorème de la bijection continue, elle réalise une bijection de \(]0,+\infty[\) sur \(\R\).
}

\correction{Oral Mines -- Une fonction réciproque implicite}{
La fonction \(f(x)=x^3+x\) est continue, strictement croissante et vérifie
\[
\lim_{x\to-\infty}f(x)=-\infty,\qquad
\lim_{x\to+\infty}f(x)=+\infty.
\]
Elle est donc bijective de \(\R\) sur \(\R\). Sa réciproque \(g\) est définie sur \(\R\).
Comme \(f(0)=0\), on a
\[
g(0)=0.
\]
La fonction \(f\) est impaire. Si \(y=f(x)\), alors \(-y=f(-x)\). Donc
\[
g(-y)=-g(y),
\]
et \(g\) est impaire.
Enfin,
\[
g(y)=1
\Longleftrightarrow
y=f(1)=2.
\]
}

\correction{Oral Centrale -- Paramètre et TVI}{
Pour \(a\in\R\), posons
\[
f(x)=x^3+x.
\]
D'après l'exercice précédent, \(f\) est une bijection de \(\R\) sur \(\R\). Donc, pour tout \(a\in\R\), il existe un unique \(x\in\R\) tel que
\[
f(x)=a,
\]
c'est-à-dire
\[
x^3+x=a.
\]
On définit alors
\[
\varphi(a)=\text{l'unique solution de }x^3+x=a.
\]
La fonction \(\varphi\) est précisément
\[
\varphi=f^{-1}.
\]
}

\correction{Oral Mines -- Fonction définie par morceaux}{
Sur \(\R^*\), la fonction \(x\mapsto\sin(1/x)\) est composée de fonctions continues, donc continue. Elle n'est pas continue en \(0\), car elle n'a pas de limite en \(0\).
Pour
\[
g(x)=x\sin(1/x),\quad g(0)=0,
\]
on a
\[
-|x|\leq g(x)\leq |x|.
\]
Donc \(g(x)\to0=g(0)\). Ainsi \(g\) est continue en \(0\).
}

\correction{Oral Centrale -- Prolongement et paramètre}{
On a déjà montré que
\[
\frac{\sqrt{1+x}-1}{x}
=
\frac1{\sqrt{1+x}+1}
\]
pour \(x\neq0\). Donc
\[
\lim_{x\to0}f(x)=\frac12.
\]
Il faut donc choisir
\[
a=\frac12.
\]
Pour cette valeur, \(f\) est continue en \(0\). La fonction \(g(x)=bf(x)+x\) est alors continue en \(0\) comme combinaison linéaire de fonctions continues. Enfin
\[
g(0)=b f(0)+0=\frac b2.
\]
}

\correction{Oral CCP -- Image d'un intervalle}{
La fonction
\[
f(x)=x^2-2x+3=(x-1)^2+2
\]
est polynomiale, donc continue. Sur \([0,2]\), son minimum vaut \(2\), atteint en \(1\). Aux bornes,
\[
f(0)=3,\qquad f(2)=3.
\]
Donc le maximum vaut \(3\). Ainsi
\[
f([0,2])=[2,3].
\]
C'est un segment, conformément au théorème de l'image continue d'un segment.
}

\correction{Oral Mines -- Non-existence par suites}{
Prenons
\[
x_n=\frac1{2n\pi},
\qquad
y_n=\frac1{\pi+2n\pi}.
\]
Alors \(x_n\to0\) et \(y_n\to0\). Mais
\[
\cos(1/x_n)=\cos(2n\pi)=1,
\]
tandis que
\[
\cos(1/y_n)=\cos(\pi+2n\pi)=-1.
\]
Les deux limites obtenues sont distinctes. Donc la limite en \(0\) n'existe pas.
}

\correction{Oral Centrale -- Bijection trigonométrique}{
La fonction sinus est continue sur \(\R\), donc sur
\[
\left[-\frac{\pi}{2},\frac{\pi}{2}\right].
\]
Elle est strictement croissante sur cet intervalle. De plus,
\[
\sin\left(-\frac{\pi}{2}\right)=-1,
\qquad
\sin\left(\frac{\pi}{2}\right)=1.
\]
Donc son image est \([-1,1]\). Par le théorème de la bijection continue, elle est bijective de
\[
\left[-\frac{\pi}{2},\frac{\pi}{2}\right]
\]
sur \([-1,1]\). Sa réciproque est la fonction
\[
\arcsin:[-1,1]\to\left[-\frac{\pi}{2},\frac{\pi}{2}\right].
\]
}

\correction{Oral Mines -- Fonction exponentielle}{
La fonction exponentielle est continue et strictement croissante sur \(\R\). De plus,
\[
\lim_{x\to-\infty}e^x=0,
\qquad
\lim_{x\to+\infty}e^x=+\infty.
\]
Son image est donc \(]0,+\infty[\). Par le théorème de la bijection continue,
\[
\exp:\R\to]0,+\infty[
\]
est bijective. Sa réciproque est
\[
\ln:]0,+\infty[\to\R.
\]
Elle est continue et strictement croissante.
}

\correction{Problème long 1 -- Étude complète d'une bijection}{
La fonction \(f(x)=x^3+x\) est polynomiale, donc continue. Pour \(x<y\),
\[
f(y)-f(x)=y^3-x^3+y-x
=(y-x)(y^2+xy+x^2+1).
\]
Or
\[
y^2+xy+x^2=\frac12(x+y)^2+\frac12(x^2+y^2)\geq0.
\]
Donc \(f(y)-f(x)>0\). La fonction est strictement croissante.

Les limites sont
\[
\lim_{x\to+\infty}f(x)=+\infty,
\qquad
\lim_{x\to-\infty}f(x)=-\infty.
\]
Par le théorème de la bijection continue, \(f\) est une bijection de \(\R\) sur \(\R\). Sa réciproque \(g\) est continue et strictement croissante.

Comme \(f\) est impaire, sa réciproque est impaire. Enfin
\[
g(y)=2
\Longleftrightarrow
y=f(2)=8+2=10.
\]
}

\correction{Problème long 2 -- TVI, unicité et paramètre}{
Soit
\[
f(x)=x^5+x.
\]
Elle est continue sur \(\R\), avec
\[
\lim_{x\to-\infty}f(x)=-\infty,
\qquad
\lim_{x\to+\infty}f(x)=+\infty.
\]
Pour tout \(a\in\R\), le TVI donne au moins une solution à \(f(x)=a\).

Si \(x<y\), alors
\[
f(y)-f(x)=y^5-x^5+y-x.
\]
La fonction \(x\mapsto x^5\) est strictement croissante, et \(x\mapsto x\) aussi. Donc
\[
f(y)>f(x).
\]
Ainsi \(f\) est strictement croissante, donc injective. La solution est unique.

On définit
\[
\varphi(a)=\text{l'unique réel }x\text{ tel que }x^5+x=a.
\]
Alors
\[
\varphi=f^{-1}.
\]
Comme \(f\) est continue et strictement croissante de \(\R\) sur \(\R\), sa réciproque \(\varphi\) est continue et strictement croissante.
}

\correction{Problème long 3 -- Prolongement et continuité}{
Le domaine de définition est
\[
]-1,+\infty[\setminus\{0\}.
\]
Pour \(x\neq0\),
\[
f(x)=\frac{\sqrt{1+x}-1}{x}
\cdot
\frac{\sqrt{1+x}+1}{\sqrt{1+x}+1}
=
\frac1{\sqrt{1+x}+1}.
\]
Donc
\[
\lim_{x\to0}f(x)=\frac12.
\]
On définit
\[
\widetilde f(x)=f(x)\quad(x\neq0),
\qquad
\widetilde f(0)=\frac12.
\]
Ainsi \(\widetilde f\) est continue en \(0\). Elle est aussi continue sur \(]-1,0[\) et \(]0,+\infty[\) comme composée et quotient de fonctions continues. Donc elle est continue sur \(]-1,+\infty[\).
}

\correction{Problème long 4 -- Fonction réciproque sur un intervalle}{
La fonction \(f(x)=x^2+1\) est continue sur \([0,+\infty[\). Elle est strictement croissante sur cet intervalle. En effet, si \(0\leq x<y\), alors
\[
y^2-x^2=(y-x)(y+x)>0.
\]
De plus,
\[
f(0)=1,\qquad \lim_{x\to+\infty}f(x)=+\infty.
\]
Donc
\[
f([0,+\infty[)=[1,+\infty[.
\]
Ainsi \(f\) est bijective de \([0,+\infty[\) sur \([1,+\infty[\).

Pour \(y\geq1\), résoudre \(y=x^2+1\) avec \(x\geq0\) donne
\[
x=\sqrt{y-1}.
\]
Donc
\[
f^{-1}(y)=\sqrt{y-1}.
\]
Vérifications :
\[
f^{-1}(f(x))=\sqrt{x^2+1-1}=\sqrt{x^2}=x
\]
car \(x\geq0\), et
\[
f(f^{-1}(y))=(\sqrt{y-1})^2+1=y.
\]
}

\correction{Problème long 5 -- Limites latérales, continuité et recollement}{
Sur \(]-\infty,1[\), \(f(x)=x^2+a\), donc \(f\) est continue. Sur \(]1,+\infty[\), \(f(x)=bx+1\), donc \(f\) est continue.

À gauche de \(1\),
\[
\lim_{x\to1^-}f(x)=1+a.
\]
À droite de \(1\),
\[
\lim_{x\to1^+}f(x)=b+1.
\]
De plus
\[
f(1)=b+1.
\]
La continuité en \(1\) équivaut donc à
\[
1+a=b+1,
\]
c'est-à-dire
\[
a=b.
\]
Trois exemples :
\[
(a,b)=(0,0),\qquad (1,1),\qquad (-2,-2).
\]
}

\end{document}